%%%%%%%%%%%%%%%%%%%%%%%%%%%%% Define Article %%%%%%%%%%%%%%%%%%%%%%%%%%%%%%%%%%
\documentclass{article}
%%%%%%%%%%%%%%%%%%%%%%%%%%%%%%%%%%%%%%%%%%%%%%%%%%%%%%%%%%%%%%%%%%%%%%%%%%%%%%%

%%%%%%%%%%%%%%%%%%%%%%%%%%%%% Using Packages %%%%%%%%%%%%%%%%%%%%%%%%%%%%%%%%%%
\usepackage[table]{xcolor}
\usepackage{tabularx}
\usepackage{array}
\usepackage{geometry}
\usepackage[utf8]{inputenc}
\usepackage[T1]{fontenc}
\usepackage{amssymb}
\usepackage{amsmath}
\usepackage{graphicx}
\usepackage{fancyhdr}
\usepackage[english]{babel}
\usepackage{mdframed}
\usepackage{titlesec}
\usepackage{enumitem}
\usepackage[breaklinks]{hyperref}

%%%%%%%%%%%%%%%%%%%%%%%%%% Colors: one blue theme %%%%%%%%%%%%%%%%%%%%%%%%%%%%%
\definecolor{deepBlue}{RGB}{31,78,140}      % headings, rules, table headers
\definecolor{midBlue}{RGB}{61,124,222}      % links, accents
\definecolor{shallowBlue}{RGB}{235,243,252} % shaded rows, boxes
\definecolor{textGray}{RGB}{90,90,90}       % header/footer text
%%%%%%%%%%%%%%%%%%%%%%%%%%%%%%%%%%%%%%%%%%%%%%%%%%%%%%%%%%%%%%%%%%%%%%%%%%%%%%%

% Other Settings
\hypersetup{
    colorlinks = true,
    linkcolor = midBlue,
    urlcolor = midBlue,
}
\urlstyle{same}

%%%%%%%%%%%%%%%%%%%%%%%%%% Page Setting %%%%%%%%%%%%%%%%%%%%%%%%%%%%%%%%%%%%%%%
\geometry{a4paper, margin=2.2cm}
\setlength{\headheight}{14pt}
\pagestyle{fancy}
\fancyhead{}
\fancyhead[L]{\color{textGray}\small ROB 538 - Multiagent Systems}
\fancyhead[C]{\color{textGray}\small Project Proposal}
\fancyhead[R]{\color{textGray}\small Mushtaq \& Khawaja}
\fancyfoot{}
\fancyfoot[C]{\color{textGray}\thepage}
\renewcommand{\headrule}{{\color{deepBlue}\hrule height 0.6pt}}

\titleformat{\section}{\Large\bfseries\color{deepBlue}}{}{0em}{}[\vspace{-0.6em}\textcolor{deepBlue}{\rule{\linewidth}{0.6pt}}]
\titlespacing*{\section}{0pt}{1.4em}{0.8em}
\titleformat{\subsection}{\large\bfseries\color{deepBlue}}{}{0em}{}
\setlist{itemsep=3pt,topsep=4pt}
\setlist[itemize]{label=\textcolor{deepBlue}{\textbullet}}

%%%%%%%%%%%%%%%%%%%%%%%%%%%% Custom helpers %%%%%%%%%%%%%%%%%%%%%%%%%%%%%%%%%%%%
% In-text citation: links to the numbered entry in the Citations list, e.g. [2]
\newcommand{\citep}[1]{\ref{#1}}

% Highlighted box for the research question
\newmdenv[topline=false,bottomline=false,rightline=false,linewidth=3pt,
  linecolor=deepBlue,backgroundcolor=shallowBlue,
  innerleftmargin=14pt,innerrightmargin=14pt,innertopmargin=10pt,innerbottommargin=10pt]{rqbox}

% Tables: white bold text on a blue header row, shaded body rows
\newcommand{\tablestyle}{\renewcommand{\arraystretch}{1.35}\rowcolors{2}{shallowBlue}{white}\arrayrulecolor{deepBlue}}
\newcommand{\headrow}{\rowcolor{deepBlue}}
\newcommand{\hd}[1]{\textcolor{white}{\textbf{#1}}}
%%%%%%%%%%%%%%%%%%%%%%%%%%%%%%%%%%%%%%%%%%%%%%%%%%%%%%%%%%%%%%%%%%%%%%%%%%%%%%%

\begin{document}

%%%%%%%%%%%%%%%%%%%%%%%%%%%%%%%% Title page %%%%%%%%%%%%%%%%%%%%%%%%%%%%%%%%%%%
\begin{titlepage}
\centering
{\color{deepBlue}\rule{\linewidth}{2pt}}\\[6mm]
{\huge\bfseries\color{deepBlue} Oregon State University}\\[2mm]
{\Large\bfseries Multiagent Systems - ROB 538}\\[4mm]
{\color{deepBlue}\rule{\linewidth}{0.6pt}}

\vspace*{18mm}
{\LARGE\bfseries Final Project Proposal \& Initial Outline}\\[8mm]
\colorbox{shallowBlue}{\parbox{0.86\linewidth}{\centering\vspace{4mm}
  {\Large\bfseries\color{deepBlue} Adaptive CFR$^+$:}\\[2mm]
  {\Large Adapting Regret Updates to Local Regret Dynamics\\ in Imperfect-Information Games}
  \vspace{4mm}}}

\vspace*{16mm}
% Put the OSU logo next to this file as "OSU logo.png" to show it here
\IfFileExists{OSU logo.png}{\includegraphics[width=0.32\linewidth]{OSU logo.png}\\[12mm]}{}
{\Large Sadiqah Mushtaq \quad\textcolor{deepBlue}{\textbullet}\quad Umna Khawaja}\\[6mm]
{\large \textbf{Instructor:} Kagan Tumer}

\vfill
{\color{deepBlue}\rule{\linewidth}{0.6pt}}\\[2mm]
{\large Due: October 6, 2026}
\end{titlepage}

%===============================================================================
\section{Paper Title}

\textbf{Adaptive CFR$^+$: Adapting Regret Updates to Local Regret Dynamics in
Imperfect-Information Games}

%===============================================================================
\section{General Summary of Proposed Paper}

In games like poker, each agent sees only part of the game: its own cards, but not its
opponents'. A standard way to find good strategies for such games is Counterfactual Regret
Minimization (CFR), in which agents repeatedly play against each other and each one learns from
its \emph{regret}: how much better it could have done by choosing differently.
CFR$^+$~\citep{tammelin} is a faster version of CFR, and Predictive CFR$^+$~\citep{farina}
speeds it up further by having each agent guess its next regret and act on that guess early.

Our project will build on CFR$^+$ by investigating whether its regret updates
can be made adaptive. Instead of applying the same update behavior at
every information set, we will use the recent history of regret at each
information set to determine how aggressively that information set should be
updated. The goal is to make the algorithm respond more quickly when regret is
changing substantially while avoiding unnecessary updates when regret has
stabilized.

We will implement and evaluate this Adaptive CFR$^+$ method using OpenSpiel
on Kuhn Poker and Leduc Poker. We will compare it against CFR, CFR$^+$, and,
where feasible, predictive CFR methods. Our primary evaluation will measure
exploitability as a function of the number of iterations, allowing us to
determine whether adaptive updates reach approximate Nash equilibrium faster
than existing CFR variants. As a \textbf{bonus}, if time allows, we will test
whether the results hold in three-player poker.

\begin{rqbox}
\textbf{\color{deepBlue}Research Question.} Can adapting the CFR$^+$ regret-update strength according to local
regret dynamics improve convergence in imperfect-information games?
\end{rqbox}

\subsection*{Agent Specification}
\begin{itemize}
\item \textbf{System objective:} Minimize exploitability, using exploitability
as the primary measure of convergence toward a Nash equilibrium in
two-player zero-sum games.
\item \textbf{Payoff:} Game utility, such as chips won or lost.
\item \textbf{Optimization signal:} Counterfactual regret.
\item \textbf{Information state:} Private cards, public cards when applicable,
and betting history.
\item \textbf{Learning algorithm:} CFR, CFR$^+$, and our proposed Adaptive
CFR$^+$.
\item \textbf{Primary investigation:} Whether adaptive regret updates reduce
exploitability faster than standard CFR and CFR$^+$.
\item \textbf{Secondary investigation:} Whether the benefit of adaptation
differs between Kuhn Poker and the larger Leduc Poker game.
\item \textbf{Bonus investigation:} Whether the benefit of adaptation carries
over from two-player to three-player poker, if time allows.
\end{itemize}

\subsection*{Problem Environment}
\begin{center}
\tablestyle
\begin{tabularx}{0.9\linewidth}{|>{\bfseries}l|X|X|}
\hline
\headrow & \hd{Project 1: 2-player poker} & \hd{Project 2 (bonus): 3-player poker} \\ \hline
Observable                 & Partially  & Partially \\ \hline
Deterministic/Stochastic   & Stochastic environment, strategic reward
                           & Stochastic environment, strategic reward \\ \hline
Repetitive/Sequential      & Sequential & Sequential \\ \hline
Episodic                   & Episodic   & Episodic \\ \hline
Static/Dynamic             & Static     & Static \\ \hline
Discrete/Continuous        & Discrete   & Discrete \\ \hline
\end{tabularx}
\end{center}

%===============================================================================
\section{Background}

CFR~\citep{zinkevich} places a simple learner at every information set, a set of game
situations an agent cannot tell apart. In two-player zero-sum games, the average of the agents'
strategies converges to a Nash equilibrium. CFR$^+$ changes three things: negative regrets are
reset to zero, players update one at a time, and later iterations count more in the average.
These changes made CFR$^+$ fast enough to solve heads-up limit Texas hold'em. Kuhn poker
(3 cards, one betting round) and Leduc poker~\citep{southey} (6 cards, two betting rounds and a
public card) are small poker games that are standard for testing these algorithms. Both are
available in OpenSpiel~\citep{openspiel}.

We have already reproduced the main result of the CFR$^+$ paper on both games using OpenSpiel.
CFR$^+$ reached about \textbf{12$\times$} lower exploitability than CFR on Kuhn poker and about
\textbf{46$\times$} lower on Leduc poker.

Our proposed method builds on these ideas by making the regret update
adaptive at the information-set level. We will investigate whether the
algorithm can use recent regret behavior to determine when an information set
should receive a more aggressive or more conservative update. For example,
information sets with rapidly changing regret may benefit from larger updates,
while information sets whose regret has stabilized may benefit from smaller
updates.

Kuhn Poker and Leduc Poker provide small, controlled environments in which
these algorithms can be evaluated efficiently.

%===============================================================================
\section{Proposed Timeline}

\begin{center}
\tablestyle
\begin{tabularx}{0.9\linewidth}{|c|X|}
\hline
\headrow\hd{Week} & \hd{Milestone} \\ \hline
2    & Project proposal; CFR vs.\ CFR$^+$ baseline reproduced (\textcolor{deepBlue}{\textbf{done}}) \\ \hline
3--4 & Implement our own CFR, CFR$^+$ and Predictive CFR$^+$; \textbf{paper intro \& background} \\ \hline
5--7 & Implement the adaptive method; run main two-player Kuhn and Leduc experiments \\ \hline
8    & Run ablation studies and begin \textbf{paper draft} \\ \hline
9    & Final experiments and figures; bonus: three-player poker \\ \hline
10   & \textbf{Presentation} \\ \hline
11   & \textbf{Final paper} \\ \hline
\end{tabularx}
\end{center}

%===============================================================================
\section{Previous Work}

\begin{center}
\tablestyle
\begin{tabularx}{0.9\linewidth}{|>{\raggedright\arraybackslash}p{3.6cm}|X|}
\hline
\headrow\hd{Method} & \hd{Contribution} \\ \hline
CFR~\citep{zinkevich}               & Introduced counterfactual regret minimization for imperfect-information games. \\ \hline
CFR$^+$~\citep{tammelin}            & Regret matching$^+$, alternating updates and linear averaging; used to solve heads-up limit hold'em. The method we build on. \\ \hline
Discounted CFR~\citep{brown}        & Down-weights early iterations with a fixed discount to speed up convergence. \\ \hline
Predictive CFR$^+$~\citep{farina}   & Predicts the next regret and acts on it early, with the same prediction everywhere. \\ \hline
\end{tabularx}
\end{center}

%===============================================================================
\section{Citations}

\begin{enumerate}[label={[\arabic*]}]
  \item \label{zinkevich}M. Zinkevich, M. Johanson, M. Bowling, C. Piccione.
        \href{https://poker.cs.ualberta.ca/publications/NIPS07-cfr.pdf}{Regret Minimization
        in Games with Incomplete Information}. NeurIPS, 2007.
  \item \label{tammelin}O. Tammelin, N. Burch, M. Johanson, M. Bowling.
        \href{https://webdocs.cs.ualberta.ca/~games/poker/publications/2015-ijcai-cfrplus.pdf}{Solving
        Heads-up Limit Texas Hold'em}. IJCAI, 2015.
  \item \label{brown}N. Brown, T. Sandholm.
        \href{https://arxiv.org/abs/1809.04040}{Solving Imperfect-Information Games via
        Discounted Regret Minimization}. AAAI, 2019.
  \item \label{farina}G. Farina, C. Kroer, T. Sandholm.
        \href{https://arxiv.org/abs/2007.14358}{Faster Game Solving via Predictive Blackwell
        Approachability: Connecting Regret Matching and Mirror Descent}. AAAI, 2021.
  \item \label{southey}F. Southey et al.
        \href{https://arxiv.org/abs/1207.1411}{Bayes' Bluff: Opponent Modelling in Poker}.
        UAI, 2005.
  \item \label{openspiel}M. Lanctot et al.
        \href{https://arxiv.org/abs/1908.09453}{OpenSpiel: A Framework for Reinforcement
        Learning in Games}. arXiv, 2019.
\end{enumerate}

\end{document}
